\documentclass[11pt]{article}

    \usepackage[breakable]{tcolorbox}
    \usepackage{parskip} % Stop auto-indenting (to mimic markdown behaviour)
    

    % Basic figure setup, for now with no caption control since it's done
    % automatically by Pandoc (which extracts ![](path) syntax from Markdown).
    \usepackage{graphicx}
    % Maintain compatibility with old templates. Remove in nbconvert 6.0
    \let\Oldincludegraphics\includegraphics
    % Ensure that by default, figures have no caption (until we provide a
    % proper Figure object with a Caption API and a way to capture that
    % in the conversion process - todo).
    \usepackage{caption}
    \DeclareCaptionFormat{nocaption}{}
    \captionsetup{format=nocaption,aboveskip=0pt,belowskip=-20pt}
    \setlength{\belowcaptionskip}{-20pt}

    \usepackage[spanish]{babel}
    \usepackage{float}
    \floatplacement{figure}{H} % forces figures to be placed at the correct location
    \usepackage{xcolor} % Allow colors to be defined
    \usepackage{enumerate} % Needed for markdown enumerations to work
    \usepackage[a4paper]{geometry} % Used to adjust the document margins
    \usepackage{amsmath} % Equations
    \usepackage{amssymb} % Equations
    \usepackage{textcomp} % defines textquotesingle
    % Hack from http://tex.stackexchange.com/a/47451/13684:
    \AtBeginDocument{%
        \def\PYZsq{\textquotesingle}% Upright quotes in Pygmentized code
    }
    \usepackage{upquote} % Upright quotes for verbatim code
    \usepackage{eurosym} % defines \euro

    \usepackage{iftex}
    \ifPDFTeX
        \usepackage[T1]{fontenc}
        \usepackage{helvet}
        \renewcommand\familydefault{\sfdefault}

        \IfFileExists{alphabeta.sty}{
              \usepackage{alphabeta}
          }{
              \usepackage[mathletters]{ucs}
              \usepackage[utf8x]{inputenc}
          }
    \else
        \usepackage{fontspec}
        \usepackage{unicode-math}
        \setmainfont{Nimbus Sans}
    \fi

    \usepackage{fancyvrb} % verbatim replacement that allows latex
    \usepackage{grffile} % extends the file name processing of package graphics
                         % to support a larger range
    \makeatletter % fix for old versions of grffile with XeLaTeX
    \@ifpackagelater{grffile}{2019/11/01}
    {
      % Do nothing on new versions
    }
    {
      \def\Gread@@xetex#1{%
        \IfFileExists{"\Gin@base".bb}%
        {\Gread@eps{\Gin@base.bb}}%
        {\Gread@@xetex@aux#1}%
      }
    }
    \makeatother
    \usepackage[Export]{adjustbox} % Used to constrain images to a maximum size
    \adjustboxset{max size={\linewidth}{\paperheight}}

    % The hyperref package gives us a pdf with properly built
    % internal navigation ('pdf bookmarks' for the table of contents,
    % internal cross-reference links, web links for URLs, etc.)
    \usepackage{hyperref}
    % The default LaTeX title has an obnoxious amount of whitespace. By default,
    % titling removes some of it. It also provides customization options.
    \usepackage{titling}
    \setlength{\droptitle}{-2em}
    \usepackage{longtable} % longtable support required by pandoc >1.10
    \usepackage{booktabs}  % table support for pandoc > 1.12.2
    \usepackage{array}     % table support for pandoc >= 2.11.3
    \usepackage{calc}      % table minipage width calculation for pandoc >= 2.11.1
    \usepackage[inline]{enumitem} % IRkernel/repr support (it uses the enumerate* environment)
    \usepackage[normalem]{ulem} % ulem is needed to support strikethroughs (\sout)
                                % normalem makes italics be italics, not underlines
    \usepackage{soul}      % strikethrough (\st) support for pandoc >= 3.0.0
    \usepackage{mathrsfs}
    

    
    % Colors for the hyperref package
    \definecolor{urlcolor}{rgb}{0,.145,.698}
    \definecolor{linkcolor}{rgb}{.71,0.21,0.01}
    \definecolor{citecolor}{rgb}{.12,.54,.11}

    % ANSI colors
    \definecolor{ansi-black}{HTML}{3E424D}
    \definecolor{ansi-black-intense}{HTML}{282C36}
    \definecolor{ansi-red}{HTML}{E75C58}
    \definecolor{ansi-red-intense}{HTML}{B22B31}
    \definecolor{ansi-green}{HTML}{00A250}
    \definecolor{ansi-green-intense}{HTML}{007427}
    \definecolor{ansi-yellow}{HTML}{DDB62B}
    \definecolor{ansi-yellow-intense}{HTML}{B27D12}
    \definecolor{ansi-blue}{HTML}{208FFB}
    \definecolor{ansi-blue-intense}{HTML}{0065CA}
    \definecolor{ansi-magenta}{HTML}{D160C4}
    \definecolor{ansi-magenta-intense}{HTML}{A03196}
    \definecolor{ansi-cyan}{HTML}{60C6C8}
    \definecolor{ansi-cyan-intense}{HTML}{258F8F}
    \definecolor{ansi-white}{HTML}{C5C1B4}
    \definecolor{ansi-white-intense}{HTML}{A1A6B2}
    \definecolor{ansi-default-inverse-fg}{HTML}{FFFFFF}
    \definecolor{ansi-default-inverse-bg}{HTML}{000000}

    % common color for the border for error outputs.
    \definecolor{outerrorbackground}{HTML}{FFDFDF}

    % commands and environments needed by pandoc snippets
    % extracted from the output of `pandoc -s`
    \providecommand{\tightlist}{%
      \setlength{\itemsep}{0pt}\setlength{\parskip}{0pt}}
    \DefineVerbatimEnvironment{Highlighting}{Verbatim}{commandchars=\\\{\}}
    % Add ',fontsize=\small' for more characters per line
    \newenvironment{Shaded}{}{}
    \newcommand{\KeywordTok}[1]{\textcolor[rgb]{0.00,0.44,0.13}{\textbf{{#1}}}}
    \newcommand{\DataTypeTok}[1]{\textcolor[rgb]{0.56,0.13,0.00}{{#1}}}
    \newcommand{\DecValTok}[1]{\textcolor[rgb]{0.25,0.63,0.44}{{#1}}}
    \newcommand{\BaseNTok}[1]{\textcolor[rgb]{0.25,0.63,0.44}{{#1}}}
    \newcommand{\FloatTok}[1]{\textcolor[rgb]{0.25,0.63,0.44}{{#1}}}
    \newcommand{\CharTok}[1]{\textcolor[rgb]{0.25,0.44,0.63}{{#1}}}
    \newcommand{\StringTok}[1]{\textcolor[rgb]{0.25,0.44,0.63}{{#1}}}
    \newcommand{\CommentTok}[1]{\textcolor[rgb]{0.38,0.63,0.69}{\textit{{#1}}}}
    \newcommand{\OtherTok}[1]{\textcolor[rgb]{0.00,0.44,0.13}{{#1}}}
    \newcommand{\AlertTok}[1]{\textcolor[rgb]{1.00,0.00,0.00}{\textbf{{#1}}}}
    \newcommand{\FunctionTok}[1]{\textcolor[rgb]{0.02,0.16,0.49}{{#1}}}
    \newcommand{\RegionMarkerTok}[1]{{#1}}
    \newcommand{\ErrorTok}[1]{\textcolor[rgb]{1.00,0.00,0.00}{\textbf{{#1}}}}
    \newcommand{\NormalTok}[1]{{#1}}

    % Additional commands for more recent versions of Pandoc
    \newcommand{\ConstantTok}[1]{\textcolor[rgb]{0.53,0.00,0.00}{{#1}}}
    \newcommand{\SpecialCharTok}[1]{\textcolor[rgb]{0.25,0.44,0.63}{{#1}}}
    \newcommand{\VerbatimStringTok}[1]{\textcolor[rgb]{0.25,0.44,0.63}{{#1}}}
    \newcommand{\SpecialStringTok}[1]{\textcolor[rgb]{0.73,0.40,0.53}{{#1}}}
    \newcommand{\ImportTok}[1]{{#1}}
    \newcommand{\DocumentationTok}[1]{\textcolor[rgb]{0.73,0.13,0.13}{\textit{{#1}}}}
    \newcommand{\AnnotationTok}[1]{\textcolor[rgb]{0.38,0.63,0.69}{\textbf{\textit{{#1}}}}}
    \newcommand{\CommentVarTok}[1]{\textcolor[rgb]{0.38,0.63,0.69}{\textbf{\textit{{#1}}}}}
    \newcommand{\VariableTok}[1]{\textcolor[rgb]{0.10,0.09,0.49}{{#1}}}
    \newcommand{\ControlFlowTok}[1]{\textcolor[rgb]{0.00,0.44,0.13}{\textbf{{#1}}}}
    \newcommand{\OperatorTok}[1]{\textcolor[rgb]{0.40,0.40,0.40}{{#1}}}
    \newcommand{\BuiltInTok}[1]{{#1}}
    \newcommand{\ExtensionTok}[1]{{#1}}
    \newcommand{\PreprocessorTok}[1]{\textcolor[rgb]{0.74,0.48,0.00}{{#1}}}
    \newcommand{\AttributeTok}[1]{\textcolor[rgb]{0.49,0.56,0.16}{{#1}}}
    \newcommand{\InformationTok}[1]{\textcolor[rgb]{0.38,0.63,0.69}{\textbf{\textit{{#1}}}}}
    \newcommand{\WarningTok}[1]{\textcolor[rgb]{0.38,0.63,0.69}{\textbf{\textit{{#1}}}}}


    % Define a nice break command that doesn't care if a line doesn't already
    % exist.
    \def\br{\hspace*{\fill} \\* }
    % Math Jax compatibility definitions
    \def\gt{>}
    \def\lt{<}
    \let\Oldtex\TeX
    \let\Oldlatex\LaTeX
    \renewcommand{\TeX}{\textrm{\Oldtex}}
    \renewcommand{\LaTeX}{\textrm{\Oldlatex}}
    % Document parameters
    % Document title
    % %%% \title{
    %     \Huge{\bfseries tpe \\ 
    %     \large{\mdseries }}
    % }

    %
    % 
    % 
    % 
    % 
    % 
    % 
    
% Pygments definitions
\makeatletter
\def\PY@reset{\let\PY@it=\relax \let\PY@bf=\relax%
    \let\PY@ul=\relax \let\PY@tc=\relax%
    \let\PY@bc=\relax \let\PY@ff=\relax}
\def\PY@tok#1{\csname PY@tok@#1\endcsname}
\def\PY@toks#1+{\ifx\relax#1\empty\else%
    \PY@tok{#1}\expandafter\PY@toks\fi}
\def\PY@do#1{\PY@bc{\PY@tc{\PY@ul{%
    \PY@it{\PY@bf{\PY@ff{#1}}}}}}}
\def\PY#1#2{\PY@reset\PY@toks#1+\relax+\PY@do{#2}}

\@namedef{PY@tok@w}{\def\PY@tc##1{\textcolor[rgb]{0.73,0.73,0.73}{##1}}}
\@namedef{PY@tok@c}{\let\PY@it=\textit\def\PY@tc##1{\textcolor[rgb]{0.24,0.48,0.48}{##1}}}
\@namedef{PY@tok@cp}{\def\PY@tc##1{\textcolor[rgb]{0.61,0.40,0.00}{##1}}}
\@namedef{PY@tok@k}{\let\PY@bf=\textbf\def\PY@tc##1{\textcolor[rgb]{0.00,0.50,0.00}{##1}}}
\@namedef{PY@tok@kp}{\def\PY@tc##1{\textcolor[rgb]{0.00,0.50,0.00}{##1}}}
\@namedef{PY@tok@kt}{\def\PY@tc##1{\textcolor[rgb]{0.69,0.00,0.25}{##1}}}
\@namedef{PY@tok@o}{\def\PY@tc##1{\textcolor[rgb]{0.40,0.40,0.40}{##1}}}
\@namedef{PY@tok@ow}{\let\PY@bf=\textbf\def\PY@tc##1{\textcolor[rgb]{0.67,0.13,1.00}{##1}}}
\@namedef{PY@tok@nb}{\def\PY@tc##1{\textcolor[rgb]{0.00,0.50,0.00}{##1}}}
\@namedef{PY@tok@nf}{\def\PY@tc##1{\textcolor[rgb]{0.00,0.00,1.00}{##1}}}
\@namedef{PY@tok@nc}{\let\PY@bf=\textbf\def\PY@tc##1{\textcolor[rgb]{0.00,0.00,1.00}{##1}}}
\@namedef{PY@tok@nn}{\let\PY@bf=\textbf\def\PY@tc##1{\textcolor[rgb]{0.00,0.00,1.00}{##1}}}
\@namedef{PY@tok@ne}{\let\PY@bf=\textbf\def\PY@tc##1{\textcolor[rgb]{0.80,0.25,0.22}{##1}}}
\@namedef{PY@tok@nv}{\def\PY@tc##1{\textcolor[rgb]{0.10,0.09,0.49}{##1}}}
\@namedef{PY@tok@no}{\def\PY@tc##1{\textcolor[rgb]{0.53,0.00,0.00}{##1}}}
\@namedef{PY@tok@nl}{\def\PY@tc##1{\textcolor[rgb]{0.46,0.46,0.00}{##1}}}
\@namedef{PY@tok@ni}{\let\PY@bf=\textbf\def\PY@tc##1{\textcolor[rgb]{0.44,0.44,0.44}{##1}}}
\@namedef{PY@tok@na}{\def\PY@tc##1{\textcolor[rgb]{0.41,0.47,0.13}{##1}}}
\@namedef{PY@tok@nt}{\let\PY@bf=\textbf\def\PY@tc##1{\textcolor[rgb]{0.00,0.50,0.00}{##1}}}
\@namedef{PY@tok@nd}{\def\PY@tc##1{\textcolor[rgb]{0.67,0.13,1.00}{##1}}}
\@namedef{PY@tok@s}{\def\PY@tc##1{\textcolor[rgb]{0.73,0.13,0.13}{##1}}}
\@namedef{PY@tok@sd}{\let\PY@it=\textit\def\PY@tc##1{\textcolor[rgb]{0.73,0.13,0.13}{##1}}}
\@namedef{PY@tok@si}{\let\PY@bf=\textbf\def\PY@tc##1{\textcolor[rgb]{0.64,0.35,0.47}{##1}}}
\@namedef{PY@tok@se}{\let\PY@bf=\textbf\def\PY@tc##1{\textcolor[rgb]{0.67,0.36,0.12}{##1}}}
\@namedef{PY@tok@sr}{\def\PY@tc##1{\textcolor[rgb]{0.64,0.35,0.47}{##1}}}
\@namedef{PY@tok@ss}{\def\PY@tc##1{\textcolor[rgb]{0.10,0.09,0.49}{##1}}}
\@namedef{PY@tok@sx}{\def\PY@tc##1{\textcolor[rgb]{0.00,0.50,0.00}{##1}}}
\@namedef{PY@tok@m}{\def\PY@tc##1{\textcolor[rgb]{0.40,0.40,0.40}{##1}}}
\@namedef{PY@tok@gh}{\let\PY@bf=\textbf\def\PY@tc##1{\textcolor[rgb]{0.00,0.00,0.50}{##1}}}
\@namedef{PY@tok@gu}{\let\PY@bf=\textbf\def\PY@tc##1{\textcolor[rgb]{0.50,0.00,0.50}{##1}}}
\@namedef{PY@tok@gd}{\def\PY@tc##1{\textcolor[rgb]{0.63,0.00,0.00}{##1}}}
\@namedef{PY@tok@gi}{\def\PY@tc##1{\textcolor[rgb]{0.00,0.52,0.00}{##1}}}
\@namedef{PY@tok@gr}{\def\PY@tc##1{\textcolor[rgb]{0.89,0.00,0.00}{##1}}}
\@namedef{PY@tok@ge}{\let\PY@it=\textit}
\@namedef{PY@tok@gs}{\let\PY@bf=\textbf}
\@namedef{PY@tok@ges}{\let\PY@bf=\textbf\let\PY@it=\textit}
\@namedef{PY@tok@gp}{\let\PY@bf=\textbf\def\PY@tc##1{\textcolor[rgb]{0.00,0.00,0.50}{##1}}}
\@namedef{PY@tok@go}{\def\PY@tc##1{\textcolor[rgb]{0.44,0.44,0.44}{##1}}}
\@namedef{PY@tok@gt}{\def\PY@tc##1{\textcolor[rgb]{0.00,0.27,0.87}{##1}}}
\@namedef{PY@tok@err}{\def\PY@bc##1{{\setlength{\fboxsep}{\string -\fboxrule}\fcolorbox[rgb]{1.00,0.00,0.00}{1,1,1}{\strut ##1}}}}
\@namedef{PY@tok@kc}{\let\PY@bf=\textbf\def\PY@tc##1{\textcolor[rgb]{0.00,0.50,0.00}{##1}}}
\@namedef{PY@tok@kd}{\let\PY@bf=\textbf\def\PY@tc##1{\textcolor[rgb]{0.00,0.50,0.00}{##1}}}
\@namedef{PY@tok@kn}{\let\PY@bf=\textbf\def\PY@tc##1{\textcolor[rgb]{0.00,0.50,0.00}{##1}}}
\@namedef{PY@tok@kr}{\let\PY@bf=\textbf\def\PY@tc##1{\textcolor[rgb]{0.00,0.50,0.00}{##1}}}
\@namedef{PY@tok@bp}{\def\PY@tc##1{\textcolor[rgb]{0.00,0.50,0.00}{##1}}}
\@namedef{PY@tok@fm}{\def\PY@tc##1{\textcolor[rgb]{0.00,0.00,1.00}{##1}}}
\@namedef{PY@tok@vc}{\def\PY@tc##1{\textcolor[rgb]{0.10,0.09,0.49}{##1}}}
\@namedef{PY@tok@vg}{\def\PY@tc##1{\textcolor[rgb]{0.10,0.09,0.49}{##1}}}
\@namedef{PY@tok@vi}{\def\PY@tc##1{\textcolor[rgb]{0.10,0.09,0.49}{##1}}}
\@namedef{PY@tok@vm}{\def\PY@tc##1{\textcolor[rgb]{0.10,0.09,0.49}{##1}}}
\@namedef{PY@tok@sa}{\def\PY@tc##1{\textcolor[rgb]{0.73,0.13,0.13}{##1}}}
\@namedef{PY@tok@sb}{\def\PY@tc##1{\textcolor[rgb]{0.73,0.13,0.13}{##1}}}
\@namedef{PY@tok@sc}{\def\PY@tc##1{\textcolor[rgb]{0.73,0.13,0.13}{##1}}}
\@namedef{PY@tok@dl}{\def\PY@tc##1{\textcolor[rgb]{0.73,0.13,0.13}{##1}}}
\@namedef{PY@tok@s2}{\def\PY@tc##1{\textcolor[rgb]{0.73,0.13,0.13}{##1}}}
\@namedef{PY@tok@sh}{\def\PY@tc##1{\textcolor[rgb]{0.73,0.13,0.13}{##1}}}
\@namedef{PY@tok@s1}{\def\PY@tc##1{\textcolor[rgb]{0.73,0.13,0.13}{##1}}}
\@namedef{PY@tok@mb}{\def\PY@tc##1{\textcolor[rgb]{0.40,0.40,0.40}{##1}}}
\@namedef{PY@tok@mf}{\def\PY@tc##1{\textcolor[rgb]{0.40,0.40,0.40}{##1}}}
\@namedef{PY@tok@mh}{\def\PY@tc##1{\textcolor[rgb]{0.40,0.40,0.40}{##1}}}
\@namedef{PY@tok@mi}{\def\PY@tc##1{\textcolor[rgb]{0.40,0.40,0.40}{##1}}}
\@namedef{PY@tok@il}{\def\PY@tc##1{\textcolor[rgb]{0.40,0.40,0.40}{##1}}}
\@namedef{PY@tok@mo}{\def\PY@tc##1{\textcolor[rgb]{0.40,0.40,0.40}{##1}}}
\@namedef{PY@tok@ch}{\let\PY@it=\textit\def\PY@tc##1{\textcolor[rgb]{0.24,0.48,0.48}{##1}}}
\@namedef{PY@tok@cm}{\let\PY@it=\textit\def\PY@tc##1{\textcolor[rgb]{0.24,0.48,0.48}{##1}}}
\@namedef{PY@tok@cpf}{\let\PY@it=\textit\def\PY@tc##1{\textcolor[rgb]{0.24,0.48,0.48}{##1}}}
\@namedef{PY@tok@c1}{\let\PY@it=\textit\def\PY@tc##1{\textcolor[rgb]{0.24,0.48,0.48}{##1}}}
\@namedef{PY@tok@cs}{\let\PY@it=\textit\def\PY@tc##1{\textcolor[rgb]{0.24,0.48,0.48}{##1}}}

\def\PYZbs{\char`\\}
\def\PYZus{\char`\_}
\def\PYZob{\char`\{}
\def\PYZcb{\char`\}}
\def\PYZca{\char`\^}
\def\PYZam{\char`\&}
\def\PYZlt{\char`\<}
\def\PYZgt{\char`\>}
\def\PYZsh{\char`\#}
\def\PYZpc{\char`\%}
\def\PYZdl{\char`\$}
\def\PYZhy{\char`\-}
\def\PYZsq{\char`\'}
\def\PYZdq{\char`\"}
\def\PYZti{\char`\~}
% for compatibility with earlier versions
\def\PYZat{@}
\def\PYZlb{[}
\def\PYZrb{]}
\makeatother


    % For linebreaks inside Verbatim environment from package fancyvrb.
    \makeatletter
        \newbox\Wrappedcontinuationbox
        \newbox\Wrappedvisiblespacebox
        \newcommand*\Wrappedvisiblespace {\textcolor{red}{\textvisiblespace}}
        \newcommand*\Wrappedcontinuationsymbol {\textcolor{red}{\llap{\tiny$\m@th\hookrightarrow$}}}
        \newcommand*\Wrappedcontinuationindent {3ex }
        \newcommand*\Wrappedafterbreak {\kern\Wrappedcontinuationindent\copy\Wrappedcontinuationbox}
        % Take advantage of the already applied Pygments mark-up to insert
        % potential linebreaks for TeX processing.
        %        {, <, #, %, $, ' and ": go to next line.
        %        _, }, ^, &, >, - and ~: stay at end of broken line.
        % Use of \textquotesingle for straight quote.
        \newcommand*\Wrappedbreaksatspecials {%
            \def\PYGZus{\discretionary{\char`\_}{\Wrappedafterbreak}{\char`\_}}%
            \def\PYGZob{\discretionary{}{\Wrappedafterbreak\char`\{}{\char`\{}}%
            \def\PYGZcb{\discretionary{\char`\}}{\Wrappedafterbreak}{\char`\}}}%
            \def\PYGZca{\discretionary{\char`\^}{\Wrappedafterbreak}{\char`\^}}%
            \def\PYGZam{\discretionary{\char`\&}{\Wrappedafterbreak}{\char`\&}}%
            \def\PYGZlt{\discretionary{}{\Wrappedafterbreak\char`\<}{\char`\<}}%
            \def\PYGZgt{\discretionary{\char`\>}{\Wrappedafterbreak}{\char`\>}}%
            \def\PYGZsh{\discretionary{}{\Wrappedafterbreak\char`\#}{\char`\#}}%
            \def\PYGZpc{\discretionary{}{\Wrappedafterbreak\char`\%}{\char`\%}}%
            \def\PYGZdl{\discretionary{}{\Wrappedafterbreak\char`\$}{\char`\$}}%
            \def\PYGZhy{\discretionary{\char`\-}{\Wrappedafterbreak}{\char`\-}}%
            \def\PYGZsq{\discretionary{}{\Wrappedafterbreak\textquotesingle}{\textquotesingle}}%
            \def\PYGZdq{\discretionary{}{\Wrappedafterbreak\char`\"}{\char`\"}}%
            \def\PYGZti{\discretionary{\char`\~}{\Wrappedafterbreak}{\char`\~}}%
        }
        % Some characters . , ; ? ! / are not pygmentized.
        % This macro makes them "active" and they will insert potential linebreaks
        \newcommand*\Wrappedbreaksatpunct {%
            \lccode`\~`\.\lowercase{\def~}{\discretionary{\hbox{\char`\.}}{\Wrappedafterbreak}{\hbox{\char`\.}}}%
            \lccode`\~`\,\lowercase{\def~}{\discretionary{\hbox{\char`\,}}{\Wrappedafterbreak}{\hbox{\char`\,}}}%
            \lccode`\~`\;\lowercase{\def~}{\discretionary{\hbox{\char`\;}}{\Wrappedafterbreak}{\hbox{\char`\;}}}%
            \lccode`\~`\:\lowercase{\def~}{\discretionary{\hbox{\char`\:}}{\Wrappedafterbreak}{\hbox{\char`\:}}}%
            \lccode`\~`\?\lowercase{\def~}{\discretionary{\hbox{\char`\?}}{\Wrappedafterbreak}{\hbox{\char`\?}}}%
            \lccode`\~`\!\lowercase{\def~}{\discretionary{\hbox{\char`\!}}{\Wrappedafterbreak}{\hbox{\char`\!}}}%
            \lccode`\~`\/\lowercase{\def~}{\discretionary{\hbox{\char`\/}}{\Wrappedafterbreak}{\hbox{\char`\/}}}%
            \catcode`\.\active
            \catcode`\,\active
            \catcode`\;\active
            \catcode`\:\active
            \catcode`\?\active
            \catcode`\!\active
            \catcode`\/\active
            \lccode`\~`\~
        }
    \makeatother

    \let\OriginalVerbatim=\Verbatim
    \makeatletter
    \renewcommand{\Verbatim}[1][1]{%
        %\parskip\z@skip
        \sbox\Wrappedcontinuationbox {\Wrappedcontinuationsymbol}%
        \sbox\Wrappedvisiblespacebox {\FV@SetupFont\Wrappedvisiblespace}%
        \def\FancyVerbFormatLine ##1{\hsize\linewidth
            \vtop{\raggedright\hyphenpenalty\z@\exhyphenpenalty\z@
                \doublehyphendemerits\z@\finalhyphendemerits\z@
                \strut ##1\strut}%
        }%
        % If the linebreak is at a space, the latter will be displayed as visible
        % space at end of first line, and a continuation symbol starts next line.
        % Stretch/shrink are however usually zero for typewriter font.
        \def\FV@Space {%
            \nobreak\hskip\z@ plus\fontdimen3\font minus\fontdimen4\font
            \discretionary{\copy\Wrappedvisiblespacebox}{\Wrappedafterbreak}
            {\kern\fontdimen2\font}%
        }%

        % Allow breaks at special characters using \PYG... macros.
        \Wrappedbreaksatspecials
        % Breaks at punctuation characters . , ; ? ! and / need catcode=\active
        \OriginalVerbatim[#1,codes*=\Wrappedbreaksatpunct]%
    }
    \makeatother

    % Exact colors from NB
    \definecolor{incolor}{HTML}{303F9F}
    \definecolor{outcolor}{HTML}{D84315}
    \definecolor{cellborder}{HTML}{CFCFCF}
    \definecolor{cellbackground}{HTML}{F7F7F7}

    % prompt
    \makeatletter
    \newcommand{\boxspacing}{\kern\kvtcb@left@rule\kern\kvtcb@boxsep}
    \makeatother
    \newcommand{\prompt}[4]{
        {\ttfamily\llap{{\color{#2}[#3]:\hspace{3pt}#4}}\vspace{-\baselineskip}}
    }
    

    
    % Prevent overflowing lines due to hard-to-break entities
    \sloppy
    % Setup hyperref package
    \hypersetup{
      breaklinks=true,  % so long urls are correctly broken across lines
      colorlinks=true,
      urlcolor=urlcolor,
      linkcolor=linkcolor,
      citecolor=citecolor,
      }
    % Slightly bigger margins than the latex defaults
    
    \geometry{
        verbose,
        tmargin=2.50cm, 
        bmargin=1.25cm,
        lmargin=1.90cm,
        rmargin=1.90cm,
        includefoot,
        heightrounded,
        footskip=1.20cm
    }

    
    

\begin{document}
    
    
    % \maketitle
    
    
    \vspace{2em}
    

    
    

    \hypertarget{trabajo-pruxe1ctico-especial}{%
\section{Trabajo Práctico Especial}\label{trabajo-pruxe1ctico-especial}}

Señales y Sistemas (TB065) - Curso 2 - Grupo 11 - 1C2025 - FIUBA\\
Martin Klöckner - mklockner@fi.uba.ar\\
Pablo Martinez Madero - pmartinezm@fi.uba.ar\\
Ernesto Dei Castelli - edei@fi.uba.ar

    \hypertarget{introducciuxf3n}{%
\subsubsection{Introducción}\label{introducciuxf3n}}

El habla humana puede ser entendida y analizada mediante modelos de
sistemas lineales e invariantes en el tiempo (sistemas LTI). En este
trabajo se considera la señal de voz como la salida de un sistema. La
respuesta de este sistema varía en función de la fuente de la voz y la
``configuración'' del tracto vocal.

Cuando la señal tiene la forma de un tren de pulsos, la salida tiene un
patrón periódico lo que corresponde a sonidos vocálicos, como por
ejemplo la el sonido \texttt{{[}a{]}}. Cuando se ve una señal ruidosa,
se pueden identificar sonidos fricativos, como por ejemplo el sonido
\texttt{{[}s{]}}.

Se realiza un análisis de la señal donde se identifican estas
diferencias a partir de una visualización en el dominio del tiempo y
posteriormente, en el dominio de la frecuencia mediante el uso de la
transformada rápida de Fourier (FFT) la cual nos permite extraer la
frecuencia fundamental (f0) y los formantes.

    \hypertarget{obtenciuxf3n-y-gruxe1fico-de-muestras}{%
\subsubsection{Obtención y Gráfico de
Muestras}\label{obtenciuxf3n-y-gruxe1fico-de-muestras}}

Para comenzar, como señal de muestra se realiza la grabación de dos
muestras de voz, una rápida y una lenta pronunciando la palabra
``Picasso'', de manera tal que la segunda tenga aproximadamente el doble
de duracion que la primera, luego se procede a graficar las señales en
el dominio del tiempo para lograr identificar regiones periódicas y no
periódicas.

Para la realización de los gráficos se utiliza el lenguaje de
programación \href{https://www.python.org}{python} junto con las
librerías \href{https://docs.scipy.org/doc/scipy/index.html}{scipy},
\href{https://matplotlib.org/}{matplotlib} y
\href{https://numpy.org/doc/stable/index.html}{numpy}.

Para empezar se importan las librerías mencionadas previamente en un
script de python, incluyendo el método
\href{https://docs.python.org/3/library/urllib.request.html\#urllib.request.urlretrieve}{urlretrieve}
del módulo
\href{https://docs.python.org/3/library/urllib.request.html\#module-urllib.request}{urllib.request}
para obtener localmente los archivos a analizar.

    \begin{tcolorbox}[breakable, size=fbox, boxrule=1pt, pad at break*=1mm,colback=cellbackground, colframe=cellborder, fontupper=\footnotesize]
\prompt{In}{incolor}{1}{\boxspacing}
\begin{Verbatim}[commandchars=\\\{\}]
\PY{k+kn}{import} \PY{n+nn}{matplotlib}\PY{n+nn}{.}\PY{n+nn}{pyplot} \PY{k}{as} \PY{n+nn}{plt}
\PY{k+kn}{import} \PY{n+nn}{numpy} \PY{k}{as} \PY{n+nn}{np}

\PY{k+kn}{from} \PY{n+nn}{scipy}\PY{n+nn}{.}\PY{n+nn}{io} \PY{k+kn}{import} \PY{n}{wavfile}
\PY{k+kn}{from} \PY{n+nn}{urllib}\PY{n+nn}{.}\PY{n+nn}{request} \PY{k+kn}{import} \PY{n}{urlretrieve}
\end{Verbatim}
\end{tcolorbox}

    Para obtener los archivos a analizar, se define una variable para
almacenar el nombre de los archivos y se descargan utilizando el método
\href{https://docs.python.org/3/library/urllib.request.html\#urllib.request.urlretrieve}{urlretrieve}
proporcionando las respectivas urls y el nombre final que se desea que
tengan los archivos una vez descargados.

    \begin{tcolorbox}[breakable, size=fbox, boxrule=1pt, pad at break*=1mm,colback=cellbackground, colframe=cellborder, fontupper=\footnotesize]
\prompt{In}{incolor}{2}{\boxspacing}
\begin{Verbatim}[commandchars=\\\{\}]
\PY{n}{short\PYZus{}sample\PYZus{}file\PYZus{}name} \PY{o}{=} \PY{l+s+s1}{\PYZsq{}}\PY{l+s+s1}{picasso\PYZus{}short.wav}\PY{l+s+s1}{\PYZsq{}}
\PY{n}{long\PYZus{}sample\PYZus{}file\PYZus{}name} \PY{o}{=} \PY{l+s+s1}{\PYZsq{}}\PY{l+s+s1}{picasso\PYZus{}long.wav}\PY{l+s+s1}{\PYZsq{}}

\PY{n}{urlretrieve}\PY{p}{(}\PY{l+s+s1}{\PYZsq{}}\PY{l+s+s1}{https://github.com/mjkloeckner/TB065/raw/main/data/picasso\PYZus{}short.wav}\PY{l+s+s1}{\PYZsq{}}\PY{p}{,} \PY{n}{short\PYZus{}sample\PYZus{}file\PYZus{}name}\PY{p}{)}\PY{p}{;}
\PY{n}{urlretrieve}\PY{p}{(}\PY{l+s+s1}{\PYZsq{}}\PY{l+s+s1}{https://github.com/mjkloeckner/TB065/raw/main/data/picasso\PYZus{}long.wav}\PY{l+s+s1}{\PYZsq{}}\PY{p}{,} \PY{n}{long\PYZus{}sample\PYZus{}file\PYZus{}name}\PY{p}{)}\PY{p}{;}
\end{Verbatim}
\end{tcolorbox}

    Continuamos leyendo el contenido del primer archivo de nombre
\texttt{picasso\_short.wav}, para lo cual se utiliza el método
\href{https://docs.scipy.org/doc/scipy/reference/generated/scipy.io.wavfile.read.html\#read}{wavfile.read}
de la librería
\href{https://docs.scipy.org/doc/scipy/index.html}{scipy}. Esta función
devuelve los datos y la tasa de muestreo leidos del archivo en formato
WAV cuyo nombre recibe como argumento.

    \begin{tcolorbox}[breakable, size=fbox, boxrule=1pt, pad at break*=1mm,colback=cellbackground, colframe=cellborder, fontupper=\footnotesize]
\prompt{In}{incolor}{3}{\boxspacing}
\begin{Verbatim}[commandchars=\\\{\}]
\PY{n}{short\PYZus{}sample\PYZus{}fs}\PY{p}{,} \PY{n}{short\PYZus{}sample\PYZus{}data} \PY{o}{=} \PY{n}{wavfile}\PY{o}{.}\PY{n}{read}\PY{p}{(}\PY{n}{short\PYZus{}sample\PYZus{}file\PYZus{}name}\PY{p}{)}
\end{Verbatim}
\end{tcolorbox}

    Para graficar los datos obtenidos resulta conveniente definir la función
\texttt{graph\_data}, la cual utiliza métodos del módulo
\href{https://matplotlib.org/stable/api/pyplot_summary.html\#module-matplotlib.pyplot}{pyplot}
de la librería \href{https://matplotlib.org/}{matplotlib} para generar
los gráficos pedidos. Esto además permite evitar patrones repetitivos en
el código,

    \begin{tcolorbox}[breakable, size=fbox, boxrule=1pt, pad at break*=1mm,colback=cellbackground, colframe=cellborder, fontupper=\footnotesize]
\prompt{In}{incolor}{4}{\boxspacing}
\begin{Verbatim}[commandchars=\\\{\}]
\PY{k}{def} \PY{n+nf}{graph\PYZus{}data}\PY{p}{(}\PY{n}{x}\PY{p}{,} \PY{n}{y}\PY{p}{,} \PY{n}{title}\PY{p}{,} \PY{n}{figure\PYZus{}number}\PY{o}{=}\PY{l+m+mi}{0}\PY{p}{,} \PY{n}{t}\PY{o}{=}\PY{l+m+mi}{0}\PY{p}{,} \PY{n}{dt}\PY{o}{=}\PY{l+m+mi}{0}\PY{p}{,} \PY{n}{a}\PY{o}{=}\PY{l+m+mi}{0}\PY{p}{,} \PY{n}{da}\PY{o}{=}\PY{l+m+mi}{0}\PY{p}{,} \PY{n}{stem}\PY{o}{=}\PY{k+kc}{False}\PY{p}{,} \PY{n}{tick}\PY{o}{=}\PY{l+m+mi}{0}\PY{p}{)}\PY{p}{:}
  \PY{n}{figure}\PY{p}{,} \PY{n}{axis} \PY{o}{=} \PY{n}{plt}\PY{o}{.}\PY{n}{subplots}\PY{p}{(}\PY{n}{num}\PY{o}{=}\PY{n}{title}\PY{p}{,} \PY{n}{figsize}\PY{o}{=}\PY{p}{(}\PY{l+m+mi}{12}\PY{p}{,} \PY{l+m+mi}{6}\PY{p}{)}\PY{p}{)}
  \PY{n}{figure\PYZus{}caption} \PY{o}{=} \PY{l+s+s1}{\PYZsq{}}\PY{l+s+s1}{Figura }\PY{l+s+s1}{\PYZsq{}} \PY{o}{+} \PY{n+nb}{str}\PY{p}{(}\PY{n}{figure\PYZus{}number}\PY{p}{)} \PY{o}{+} \PY{l+s+s1}{\PYZsq{}}\PY{l+s+s1}{: }\PY{l+s+s1}{\PYZsq{}} \PY{o}{+} \PY{n}{title}
  \PY{n}{figure}\PY{o}{.}\PY{n}{text}\PY{p}{(}\PY{l+m+mf}{.5}\PY{p}{,} \PY{o}{\PYZhy{}}\PY{l+m+mf}{0.02}\PY{p}{,} \PY{n}{figure\PYZus{}caption}\PY{p}{,} \PY{n}{ha}\PY{o}{=}\PY{l+s+s1}{\PYZsq{}}\PY{l+s+s1}{center}\PY{l+s+s1}{\PYZsq{}}\PY{p}{,} \PY{n}{fontsize}\PY{o}{=}\PY{l+m+mi}{12}\PY{p}{)}

  \PY{k}{if} \PY{n}{stem} \PY{o}{==} \PY{k+kc}{True}\PY{p}{:}
    \PY{n}{plt}\PY{o}{.}\PY{n}{stem}\PY{p}{(}\PY{n}{x}\PY{p}{,} \PY{n}{y}\PY{p}{,} \PY{n}{label}\PY{o}{=}\PY{l+s+s1}{\PYZsq{}}\PY{l+s+s1}{Coeficientes de la Serie de Fourier}\PY{l+s+s1}{\PYZsq{}}\PY{p}{,} \PY{n}{basefmt}\PY{o}{=}\PY{l+s+s2}{\PYZdq{}}\PY{l+s+s2}{ }\PY{l+s+s2}{\PYZdq{}}\PY{p}{)}
  \PY{k}{else}\PY{p}{:}
    \PY{n}{plt}\PY{o}{.}\PY{n}{plot}\PY{p}{(}\PY{n}{x}\PY{p}{,} \PY{n}{y}\PY{p}{,} \PY{n}{label}\PY{o}{=}\PY{l+s+s1}{\PYZsq{}}\PY{l+s+s1}{Señal de Audio}\PY{l+s+s1}{\PYZsq{}}\PY{p}{)}

  \PY{n}{axis}\PY{o}{.}\PY{n}{set}\PY{p}{(}\PY{n}{xlabel}\PY{o}{=}\PY{l+s+s1}{\PYZsq{}}\PY{l+s+s1}{Tiempo [s]}\PY{l+s+s1}{\PYZsq{}} \PY{k}{if} \PY{n}{stem} \PY{o}{==} \PY{k+kc}{False} \PY{k}{else} \PY{l+s+s2}{\PYZdq{}}\PY{l+s+s2}{Frecuencia [Hz]}\PY{l+s+s2}{\PYZdq{}}\PY{p}{,} \PY{n}{ylabel}\PY{o}{=}\PY{l+s+s1}{\PYZsq{}}\PY{l+s+s1}{Amplitud}\PY{l+s+s1}{\PYZsq{}}\PY{p}{)}
  \PY{n}{plt}\PY{o}{.}\PY{n}{grid}\PY{p}{(}\PY{k+kc}{True}\PY{p}{)}
  \PY{n}{plt}\PY{o}{.}\PY{n}{legend}\PY{p}{(}\PY{p}{)}
  \PY{n}{plt}\PY{o}{.}\PY{n}{xlim}\PY{p}{(}\PY{p}{[}\PY{n}{t}\PY{p}{,} \PY{n}{t}\PY{o}{+}\PY{n}{dt} \PY{k}{if} \PY{n}{t} \PY{o}{!=} \PY{l+m+mi}{0} \PY{k}{else} \PY{n}{x}\PY{p}{[}\PY{o}{\PYZhy{}}\PY{l+m+mi}{1}\PY{p}{]}\PY{p}{]}\PY{p}{)}
  \PY{n}{axis}\PY{o}{.}\PY{n}{axvspan}\PY{p}{(}\PY{n}{a}\PY{p}{,} \PY{n}{a}\PY{o}{+}\PY{n}{da}\PY{p}{,} \PY{n}{color}\PY{o}{=}\PY{l+s+s1}{\PYZsq{}}\PY{l+s+s1}{skyblue}\PY{l+s+s1}{\PYZsq{}}\PY{p}{,} \PY{n}{alpha}\PY{o}{=}\PY{l+m+mi}{0} \PY{k}{if} \PY{n}{a} \PY{o}{==} \PY{l+m+mi}{0} \PY{k}{else} \PY{l+m+mf}{0.50}\PY{p}{)} \PY{c+c1}{\PYZsh{} resaltado}

  \PY{k}{if} \PY{n}{tick} \PY{o}{!=} \PY{l+m+mi}{0}\PY{p}{:}
    \PY{k+kn}{from} \PY{n+nn}{matplotlib}\PY{n+nn}{.}\PY{n+nn}{ticker} \PY{k+kn}{import} \PY{n}{MultipleLocator}
    \PY{n}{plt}\PY{o}{.}\PY{n}{gca}\PY{p}{(}\PY{p}{)}\PY{o}{.}\PY{n}{xaxis}\PY{o}{.}\PY{n}{set\PYZus{}major\PYZus{}locator}\PY{p}{(}\PY{n}{MultipleLocator}\PY{p}{(}\PY{n}{tick}\PY{p}{)}\PY{p}{)}

  \PY{n}{plt}\PY{o}{.}\PY{n}{show}\PY{p}{(}\PY{p}{)}
\end{Verbatim}
\end{tcolorbox}

    Para poder graficar los datos se necesita saber a qué tiempo corresponde
cada valor de la señal; para eso, se utiliza el método
\href{https://numpy.org/doc/stable/reference/generated/numpy.arange.html\#numpy-arange}{arange}
de la libreria \href{https://numpy.org/doc/stable/index.html}{numpy}.
Este método devuelve un arreglo de números equidistantes del largo que
recibe como parámetro, en este caso del largo de los datos leídos del
archivo. Para convertirlo a segundos se divide por la tasa de muestreo,
la cual también se obtiene cuando se lee el archivo.

    \begin{tcolorbox}[breakable, size=fbox, boxrule=1pt, pad at break*=1mm,colback=cellbackground, colframe=cellborder, fontupper=\footnotesize]
\prompt{In}{incolor}{5}{\boxspacing}
\begin{Verbatim}[commandchars=\\\{\}]
\PY{n}{short\PYZus{}sample\PYZus{}time} \PY{o}{=} \PY{n}{np}\PY{o}{.}\PY{n}{arange}\PY{p}{(}\PY{n+nb}{len}\PY{p}{(}\PY{n}{short\PYZus{}sample\PYZus{}data}\PY{p}{)}\PY{p}{)} \PY{o}{/} \PY{n}{short\PYZus{}sample\PYZus{}fs}
\end{Verbatim}
\end{tcolorbox}

    Se define una variable para almacenar el título de la figura, esto para
posteriormente pasarlo como argumento a la función \texttt{graph\_data}.

    \begin{tcolorbox}[breakable, size=fbox, boxrule=1pt, pad at break*=1mm,colback=cellbackground, colframe=cellborder, fontupper=\footnotesize]
\prompt{In}{incolor}{6}{\boxspacing}
\begin{Verbatim}[commandchars=\\\{\}]
\PY{n}{title} \PY{o}{=} \PY{l+s+s1}{\PYZsq{}}\PY{l+s+s1}{Gráfico de `}\PY{l+s+s1}{\PYZsq{}} \PY{o}{+} \PY{n+nb}{str}\PY{p}{(}\PY{n}{short\PYZus{}sample\PYZus{}file\PYZus{}name}\PY{p}{)} \PY{o}{+} \PY{l+s+s1}{\PYZsq{}}\PY{l+s+s1}{` en dominio de tiempo}\PY{l+s+s1}{\PYZsq{}}
\end{Verbatim}
\end{tcolorbox}

    Finalmente se grafica el pimer archivo utilizando la funcion
\texttt{graph\_data}

    \begin{tcolorbox}[breakable, size=fbox, boxrule=1pt, pad at break*=1mm,colback=cellbackground, colframe=cellborder, fontupper=\footnotesize]
\prompt{In}{incolor}{7}{\boxspacing}
\begin{Verbatim}[commandchars=\\\{\}]
\PY{n}{graph\PYZus{}data}\PY{p}{(}\PY{n}{short\PYZus{}sample\PYZus{}time}\PY{p}{,} \PY{n}{short\PYZus{}sample\PYZus{}data}\PY{p}{,} \PY{n}{title}\PY{p}{,} \PY{l+m+mi}{1}\PY{p}{)}
\end{Verbatim}
\end{tcolorbox}

    \vspace{-0.50em}
    \begin{center}
    \adjustimage{max size={\linewidth}{\paperheight}}{tpe_files/tpe_16_0.png}
    \end{center}
    { \hspace*{\fill} \\}
    \vspace{-2.75em}
    
    Ahora realizamos exactamente el mismo procedimiento para el segundo
archivo de aproximadamente el doble de duración.

    \begin{tcolorbox}[breakable, size=fbox, boxrule=1pt, pad at break*=1mm,colback=cellbackground, colframe=cellborder, fontupper=\footnotesize]
\prompt{In}{incolor}{8}{\boxspacing}
\begin{Verbatim}[commandchars=\\\{\}]
\PY{n}{long\PYZus{}sample\PYZus{}fs}\PY{p}{,} \PY{n}{long\PYZus{}sample\PYZus{}data} \PY{o}{=} \PY{n}{wavfile}\PY{o}{.}\PY{n}{read}\PY{p}{(}\PY{n}{long\PYZus{}sample\PYZus{}file\PYZus{}name}\PY{p}{)}
\PY{n}{long\PYZus{}sample\PYZus{}time} \PY{o}{=} \PY{n}{np}\PY{o}{.}\PY{n}{arange}\PY{p}{(}\PY{n+nb}{len}\PY{p}{(}\PY{n}{long\PYZus{}sample\PYZus{}data}\PY{p}{)}\PY{p}{)} \PY{o}{/} \PY{n}{long\PYZus{}sample\PYZus{}fs}
\PY{n}{title} \PY{o}{=} \PY{l+s+s1}{\PYZsq{}}\PY{l+s+s1}{Gráfico de `}\PY{l+s+s1}{\PYZsq{}} \PY{o}{+} \PY{n+nb}{str}\PY{p}{(}\PY{n}{long\PYZus{}sample\PYZus{}file\PYZus{}name}\PY{p}{)} \PY{o}{+} \PY{l+s+s1}{\PYZsq{}}\PY{l+s+s1}{` en dominio de tiempo}\PY{l+s+s1}{\PYZsq{}}
\PY{n}{graph\PYZus{}data}\PY{p}{(}\PY{n}{long\PYZus{}sample\PYZus{}time}\PY{p}{,} \PY{n}{long\PYZus{}sample\PYZus{}data}\PY{p}{,} \PY{n}{title}\PY{p}{,} \PY{l+m+mi}{2}\PY{p}{)}
\end{Verbatim}
\end{tcolorbox}

    \vspace{-0.50em}
    \begin{center}
    \adjustimage{max size={\linewidth}{\paperheight}}{tpe_files/tpe_18_0.png}
    \end{center}
    { \hspace*{\fill} \\}
    \vspace{-2.75em}
    
    \hypertarget{analisis-visual-en-dominio-de-tiempo}{%
\subsection{Analisis Visual en Dominio de
Tiempo}\label{analisis-visual-en-dominio-de-tiempo}}

    Haciendo una inspeccion visual de las señales graficadas en las figuras
1 y 2, se identifican ciertas partes que tienen un patron que se repite
despues de un intervalo de tiempo.

Para el primer archivo, se toma como ejemplo el intervalo entre
\texttt{0.08} y \texttt{0.12} en los cuales se observa un periódo que
parece repetirse cada \texttt{0.0093s}, esto resulta en una frecuencia
de aproximadamente \texttt{107.5Hz}; dicho intervalo de la señal se
puede ver en figura 3, en la cual también se resalta el período
fundamental de la misma.

    \begin{tcolorbox}[breakable, size=fbox, boxrule=1pt, pad at break*=1mm,colback=cellbackground, colframe=cellborder, fontupper=\footnotesize]
\prompt{In}{incolor}{9}{\boxspacing}
\begin{Verbatim}[commandchars=\\\{\}]
\PY{n}{title} \PY{o}{=} \PY{l+s+s1}{\PYZsq{}}\PY{l+s+s1}{Gráfico de `}\PY{l+s+s1}{\PYZsq{}} \PY{o}{+} \PY{n+nb}{str}\PY{p}{(}\PY{n}{short\PYZus{}sample\PYZus{}file\PYZus{}name}\PY{p}{)} \PY{o}{+} \PY{l+s+s1}{\PYZsq{}}\PY{l+s+s1}{` en dominio de tiempo}\PY{l+s+s1}{\PYZsq{}}
\PY{n}{graph\PYZus{}data}\PY{p}{(}\PY{n}{short\PYZus{}sample\PYZus{}time}\PY{p}{,} \PY{n}{short\PYZus{}sample\PYZus{}data}\PY{p}{,} \PY{n}{title}\PY{p}{,} \PY{l+m+mi}{3}\PY{p}{,} \PY{n}{t}\PY{o}{=}\PY{l+m+mf}{0.08}\PY{p}{,} \PY{n}{dt}\PY{o}{=}\PY{l+m+mf}{0.04}\PY{p}{,} \PY{n}{a}\PY{o}{=}\PY{l+m+mf}{0.0935}\PY{p}{,} \PY{n}{da}\PY{o}{=}\PY{l+m+mf}{0.0073}\PY{p}{)}
\end{Verbatim}
\end{tcolorbox}

    \vspace{-0.50em}
    \begin{center}
    \adjustimage{max size={\linewidth}{\paperheight}}{tpe_files/tpe_21_0.png}
    \end{center}
    { \hspace*{\fill} \\}
    \vspace{-2.75em}
    
    De manera analoga para el segundo archivo, de mayor duración, se realiza
una inspección visual de la señal, en este caso se toma el intervalo
entre \texttt{0.08s} y \texttt{0.120s} como se observa en la figura 4,
además en la misma se grafica el periodo fundamental el cual resulta
aproximadamente \texttt{0.00756s}, resultando en una frecuencia
aproximada de \texttt{132.3Hz}

    \begin{tcolorbox}[breakable, size=fbox, boxrule=1pt, pad at break*=1mm,colback=cellbackground, colframe=cellborder, fontupper=\footnotesize]
\prompt{In}{incolor}{10}{\boxspacing}
\begin{Verbatim}[commandchars=\\\{\}]
\PY{n}{title} \PY{o}{=} \PY{l+s+s1}{\PYZsq{}}\PY{l+s+s1}{Gráfico de `}\PY{l+s+s1}{\PYZsq{}} \PY{o}{+} \PY{n+nb}{str}\PY{p}{(}\PY{n}{long\PYZus{}sample\PYZus{}file\PYZus{}name}\PY{p}{)} \PY{o}{+} \PY{l+s+s1}{\PYZsq{}}\PY{l+s+s1}{` en dominio de tiempo}\PY{l+s+s1}{\PYZsq{}}
\PY{n}{graph\PYZus{}data}\PY{p}{(}\PY{n}{long\PYZus{}sample\PYZus{}time}\PY{p}{,} \PY{n}{long\PYZus{}sample\PYZus{}data}\PY{p}{,} \PY{n}{title}\PY{p}{,} \PY{l+m+mi}{4}\PY{p}{,} \PY{n}{t}\PY{o}{=}\PY{l+m+mf}{0.08}\PY{p}{,} \PY{n}{dt}\PY{o}{=}\PY{l+m+mf}{0.04}\PY{p}{,} \PY{n}{a}\PY{o}{=}\PY{l+m+mf}{0.09525}\PY{p}{,} \PY{n}{da}\PY{o}{=}\PY{l+m+mf}{0.00756}\PY{p}{)}
\end{Verbatim}
\end{tcolorbox}

    \vspace{-0.50em}
    \begin{center}
    \adjustimage{max size={\linewidth}{\paperheight}}{tpe_files/tpe_23_0.png}
    \end{center}
    { \hspace*{\fill} \\}
    \vspace{-2.75em}
    
    Realizando un analisis similar en busca de intervalos peridicos, se
realiza una inspeccion visual de la señal, pero esta vez buscando
patrones no repetitivos.

A continuación se muestra un intervalo de la señal del archivo corto, en
donde la señal resulta no períodica, se puede ver que si bien parece
haber un patron repetitivo y asemejar a una señal senoidal, no se puede
considerar períodica porque la señal montada, o ruido, no es periodico,
por lo que la señal resultante no lo es.

    \begin{tcolorbox}[breakable, size=fbox, boxrule=1pt, pad at break*=1mm,colback=cellbackground, colframe=cellborder, fontupper=\footnotesize]
\prompt{In}{incolor}{11}{\boxspacing}
\begin{Verbatim}[commandchars=\\\{\}]
\PY{n}{title} \PY{o}{=} \PY{l+s+s1}{\PYZsq{}}\PY{l+s+s1}{Gráfico de `}\PY{l+s+s1}{\PYZsq{}} \PY{o}{+} \PY{n+nb}{str}\PY{p}{(}\PY{n}{short\PYZus{}sample\PYZus{}file\PYZus{}name}\PY{p}{)} \PY{o}{+} \PY{l+s+s1}{\PYZsq{}}\PY{l+s+s1}{` en dominio de tiempo}\PY{l+s+s1}{\PYZsq{}}
\PY{n}{graph\PYZus{}data}\PY{p}{(}\PY{n}{short\PYZus{}sample\PYZus{}time}\PY{p}{,} \PY{n}{short\PYZus{}sample\PYZus{}data}\PY{p}{,} \PY{n}{title}\PY{p}{,} \PY{l+m+mi}{5}\PY{p}{,} \PY{n}{t}\PY{o}{=}\PY{l+m+mf}{0.46}\PY{p}{,} \PY{n}{dt}\PY{o}{=}\PY{l+m+mf}{0.03}\PY{p}{)}\PY{p}{;}
\end{Verbatim}
\end{tcolorbox}

    \vspace{-0.50em}
    \begin{center}
    \adjustimage{max size={\linewidth}{\paperheight}}{tpe_files/tpe_25_0.png}
    \end{center}
    { \hspace*{\fill} \\}
    \vspace{-2.75em}
    
    Para el archivo largo se toma el intervalo entre \texttt{0.445s} y
\texttt{0.460} como ejemplo de intervalo no periodico de la señal, este
intervalo se puede ver en la figura 6 a continuación.

    \begin{tcolorbox}[breakable, size=fbox, boxrule=1pt, pad at break*=1mm,colback=cellbackground, colframe=cellborder, fontupper=\footnotesize]
\prompt{In}{incolor}{12}{\boxspacing}
\begin{Verbatim}[commandchars=\\\{\}]
\PY{n}{title} \PY{o}{=} \PY{l+s+s1}{\PYZsq{}}\PY{l+s+s1}{Gráfico de `}\PY{l+s+s1}{\PYZsq{}} \PY{o}{+} \PY{n+nb}{str}\PY{p}{(}\PY{n}{long\PYZus{}sample\PYZus{}file\PYZus{}name}\PY{p}{)} \PY{o}{+} \PY{l+s+s1}{\PYZsq{}}\PY{l+s+s1}{` en dominio de tiempo}\PY{l+s+s1}{\PYZsq{}}
\PY{n}{graph\PYZus{}data}\PY{p}{(}\PY{n}{long\PYZus{}sample\PYZus{}time}\PY{p}{,} \PY{n}{long\PYZus{}sample\PYZus{}data}\PY{p}{,} \PY{n}{title}\PY{p}{,} \PY{l+m+mi}{6}\PY{p}{,} \PY{n}{t}\PY{o}{=}\PY{l+m+mf}{0.445}\PY{p}{,} \PY{n}{dt}\PY{o}{=}\PY{l+m+mf}{0.015}\PY{p}{)}
\end{Verbatim}
\end{tcolorbox}

    \vspace{-0.50em}
    \begin{center}
    \adjustimage{max size={\linewidth}{\paperheight}}{tpe_files/tpe_27_0.png}
    \end{center}
    { \hspace*{\fill} \\}
    \vspace{-2.75em}
    
    \hypertarget{identificaciuxf3n-de-las-sonidos-vocalicos}{%
\subsubsection{Identificación de las sonidos
vocalicos}\label{identificaciuxf3n-de-las-sonidos-vocalicos}}

En la figura 7 el grafico resaltado corresponde con el sonido vocalico
\texttt{{[}a{]}} ya que se asemeja a un tren de impulsos
cuasi-periodico.

    \begin{tcolorbox}[breakable, size=fbox, boxrule=1pt, pad at break*=1mm,colback=cellbackground, colframe=cellborder, fontupper=\footnotesize]
\prompt{In}{incolor}{13}{\boxspacing}
\begin{Verbatim}[commandchars=\\\{\}]
\PY{n}{title} \PY{o}{=} \PY{l+s+s1}{\PYZsq{}}\PY{l+s+s1}{Gráfico de `}\PY{l+s+s1}{\PYZsq{}} \PY{o}{+} \PY{n+nb}{str}\PY{p}{(}\PY{n}{short\PYZus{}sample\PYZus{}file\PYZus{}name}\PY{p}{)} \PY{o}{+} \PY{l+s+s1}{\PYZsq{}}\PY{l+s+s1}{` en dominio de tiempo}\PY{l+s+s1}{\PYZsq{}}
\PY{n}{graph\PYZus{}data}\PY{p}{(}\PY{n}{short\PYZus{}sample\PYZus{}time}\PY{p}{,} \PY{n}{short\PYZus{}sample\PYZus{}data}\PY{p}{,} \PY{n}{title}\PY{p}{,} \PY{l+m+mi}{7}\PY{p}{,} \PY{n}{t}\PY{o}{=}\PY{l+m+mf}{0.35}\PY{p}{,} \PY{n}{dt}\PY{o}{=}\PY{l+m+mf}{0.05}\PY{p}{,} \PY{n}{a}\PY{o}{=}\PY{l+m+mf}{0.361}\PY{p}{,} \PY{n}{da}\PY{o}{=}\PY{l+m+mf}{0.008}\PY{p}{)}\PY{p}{;}
\end{Verbatim}
\end{tcolorbox}

    \vspace{-0.50em}
    \begin{center}
    \adjustimage{max size={\linewidth}{\paperheight}}{tpe_files/tpe_29_0.png}
    \end{center}
    { \hspace*{\fill} \\}
    \vspace{-2.75em}
    
    El sonido vocalico \texttt{{[}s{]}} corresponde con una señal de entrada
del sistema de ruido blanco, este patron se identifica en la figura 6,
en la cual se puede ver una señal senoidal la cual parece tener montada
una señal de ruidosa.

    \hypertarget{analisis-en-frecuencia}{%
\subsection{Analisis en Frecuencia}\label{analisis-en-frecuencia}}

Para calcular los coeficientes de Fourier, se utiliza los metodos
\href{https://docs.scipy.org/doc/scipy/reference/generated/scipy.fft.fft.html\#scipy.fft.fft}{fft}
y
\href{https://docs.scipy.org/doc/scipy/reference/generated/scipy.fft.fftfreq.html\#fftfreq}{fftfreq}
del modulo
\href{https://docs.scipy.org/doc/scipy/reference/fft.html\#discrete-fourier-transforms-scipy-fft}{fft}
de la libreria
\href{https://docs.scipy.org/doc/scipy/index.html}{scipy}. Se define la
funcion \texttt{graph\_fourier} la cual permite graficar tanto los
coeficientes de la serie de Fourier del intervalo definido por
\texttt{{[}x,\ y{]}} como la transformada de la misma, para el mismo
intervalo.

    \begin{tcolorbox}[breakable, size=fbox, boxrule=1pt, pad at break*=1mm,colback=cellbackground, colframe=cellborder, fontupper=\footnotesize]
\prompt{In}{incolor}{14}{\boxspacing}
\begin{Verbatim}[commandchars=\\\{\}]
\PY{k+kn}{from} \PY{n+nn}{scipy}\PY{n+nn}{.}\PY{n+nn}{fft} \PY{k+kn}{import} \PY{n}{fft}\PY{p}{,} \PY{n}{fftfreq}

\PY{k}{def} \PY{n+nf}{graph\PYZus{}fourier}\PY{p}{(}\PY{n}{data}\PY{p}{,} \PY{n}{fs}\PY{p}{,} \PY{n}{t}\PY{p}{,} \PY{n}{dt}\PY{p}{,} \PY{n}{title}\PY{p}{,} \PY{n}{fig\PYZus{}num}\PY{p}{,} \PY{n}{f\PYZus{}max}\PY{p}{,} \PY{n}{coef}\PY{o}{=}\PY{k+kc}{False}\PY{p}{,} \PY{n}{ticks}\PY{o}{=}\PY{l+m+mi}{0}\PY{p}{)}\PY{p}{:}
  \PY{n}{interval\PYZus{}data} \PY{o}{=} \PY{n}{data}\PY{p}{[}\PY{n+nb}{int}\PY{p}{(}\PY{n}{t}\PY{o}{*}\PY{n}{fs}\PY{p}{)}\PY{p}{:}\PY{n+nb}{int}\PY{p}{(}\PY{p}{(}\PY{n}{t}\PY{o}{+}\PY{n}{dt}\PY{p}{)}\PY{o}{*}\PY{n}{fs}\PY{p}{)}\PY{p}{]}
  \PY{n}{interval\PYZus{}fft} \PY{o}{=} \PY{n}{fft}\PY{p}{(}\PY{n}{interval\PYZus{}data}\PY{p}{)}
  \PY{n}{interval\PYZus{}freqs} \PY{o}{=} \PY{n}{fftfreq}\PY{p}{(}\PY{n+nb}{len}\PY{p}{(}\PY{n}{interval\PYZus{}data}\PY{p}{)}\PY{p}{,} \PY{n}{d}\PY{o}{=}\PY{l+m+mi}{1}\PY{o}{/}\PY{n}{fs}\PY{p}{)}

  \PY{c+c1}{\PYZsh{} Se toma la parte positiva en ambos casos (primer parte del arreglo)}
  \PY{n}{x} \PY{o}{=} \PY{n}{interval\PYZus{}freqs}\PY{p}{[}\PY{p}{:}\PY{n+nb}{len}\PY{p}{(}\PY{n}{interval\PYZus{}data}\PY{p}{)} \PY{o}{/}\PY{o}{/} \PY{l+m+mi}{2}\PY{p}{]}
  \PY{n}{y} \PY{o}{=} \PY{n}{np}\PY{o}{.}\PY{n}{abs}\PY{p}{(}\PY{n}{interval\PYZus{}fft}\PY{p}{[}\PY{p}{:}\PY{n+nb}{len}\PY{p}{(}\PY{n}{interval\PYZus{}data}\PY{p}{)} \PY{o}{/}\PY{o}{/} \PY{l+m+mi}{2}\PY{p}{]}\PY{p}{)}

  \PY{n}{graph\PYZus{}data}\PY{p}{(}\PY{n}{x}\PY{p}{,} \PY{n}{y}\PY{p}{,} \PY{n}{title}\PY{p}{,} \PY{n}{fig\PYZus{}num}\PY{p}{,} \PY{l+m+mf}{0.01} \PY{k}{if} \PY{n}{f\PYZus{}max} \PY{o}{!=} \PY{l+m+mi}{0} \PY{k}{else} \PY{l+m+mi}{0}\PY{p}{,} \PY{n}{f\PYZus{}max}\PY{p}{,} \PY{n}{stem}\PY{o}{=}\PY{n}{coef}\PY{p}{,} \PY{n}{tick}\PY{o}{=}\PY{n}{ticks}\PY{p}{)}
\end{Verbatim}
\end{tcolorbox}

    La palabra ``Picasso'' pronunciada en el archivo de muestra cuenta con 3
vocales: ``a'', ``i'' y ``o'', ya identificamos el intervalo en donde se
pronuncia la ``a'', el cual se muestra en la figura 7, a continuación se
utiliza la funcion \texttt{graph\_fourier} definida previamente para
calcular los coeficientes de fourier, utilizando el intervalo
\texttt{{[}0.361,\ 0.369{]}} en el cual se encuentra un solo periodo de
la señal periodica producto de pronunciar la ``a''.

    \begin{tcolorbox}[breakable, size=fbox, boxrule=1pt, pad at break*=1mm,colback=cellbackground, colframe=cellborder, fontupper=\footnotesize]
\prompt{In}{incolor}{15}{\boxspacing}
\begin{Verbatim}[commandchars=\\\{\}]
\PY{n}{t}\PY{o}{=}\PY{l+m+mf}{0.361}
\PY{n}{dt}\PY{o}{=}\PY{l+m+mf}{0.008} \PY{c+c1}{\PYZsh{} 0.361 + 0.008 = 0.369}
\PY{n}{title} \PY{o}{=} \PY{l+s+sa}{f}\PY{l+s+s1}{\PYZsq{}}\PY{l+s+s1}{Coeficientes de la Serie de Fourier para el intervalo [}\PY{l+s+si}{\PYZob{}}\PY{n}{t}\PY{l+s+si}{\PYZcb{}}\PY{l+s+s1}{, }\PY{l+s+si}{\PYZob{}}\PY{n}{t}\PY{o}{+}\PY{n}{dt}\PY{l+s+si}{\PYZcb{}}\PY{l+s+s1}{] del archivo `}\PY{l+s+si}{\PYZob{}}\PY{n}{short\PYZus{}sample\PYZus{}file\PYZus{}name}\PY{l+s+si}{\PYZcb{}}\PY{l+s+s1}{`}\PY{l+s+s1}{\PYZsq{}}
\PY{n}{graph\PYZus{}fourier}\PY{p}{(}\PY{n}{short\PYZus{}sample\PYZus{}data}\PY{p}{,} \PY{n}{short\PYZus{}sample\PYZus{}fs}\PY{p}{,} \PY{n}{t}\PY{p}{,} \PY{n}{dt}\PY{p}{,} \PY{n}{title}\PY{p}{,} \PY{l+m+mi}{8}\PY{p}{,} \PY{l+m+mi}{5000}\PY{p}{,} \PY{n}{coef}\PY{o}{=}\PY{k+kc}{True}\PY{p}{,} \PY{n}{ticks}\PY{o}{=}\PY{l+m+mi}{250}\PY{p}{)}
\end{Verbatim}
\end{tcolorbox}

    \vspace{-0.50em}
    \begin{center}
    \adjustimage{max size={\linewidth}{\paperheight}}{tpe_files/tpe_34_0.png}
    \end{center}
    { \hspace*{\fill} \\}
    \vspace{-2.75em}
    
    Tomando 4 periodos de la señal, los coeficientes de Fourier resultan
como se ve en la figura 9 a continuación, se puede ver que tiene mucho
mas ``detalle'' que la figura 7, en la cual se tomaba solo un periodo
para realizar la FFT

    \begin{tcolorbox}[breakable, size=fbox, boxrule=1pt, pad at break*=1mm,colback=cellbackground, colframe=cellborder, fontupper=\footnotesize]
\prompt{In}{incolor}{16}{\boxspacing}
\begin{Verbatim}[commandchars=\\\{\}]
\PY{n}{t}\PY{o}{=}\PY{l+m+mf}{0.3533}
\PY{n}{dt}\PY{o}{=}\PY{l+m+mf}{0.0235}
\PY{n}{title} \PY{o}{=} \PY{l+s+sa}{f}\PY{l+s+s1}{\PYZsq{}}\PY{l+s+s1}{Coeficientes de la Serie de Fourier para el intervalo [}\PY{l+s+si}{\PYZob{}}\PY{n}{t}\PY{l+s+si}{\PYZcb{}}\PY{l+s+s1}{, }\PY{l+s+si}{\PYZob{}}\PY{n}{t}\PY{o}{+}\PY{n}{dt}\PY{l+s+si}{\PYZcb{}}\PY{l+s+s1}{] del archivo `}\PY{l+s+si}{\PYZob{}}\PY{n}{short\PYZus{}sample\PYZus{}file\PYZus{}name}\PY{l+s+si}{\PYZcb{}}\PY{l+s+s1}{`}\PY{l+s+s1}{\PYZsq{}}
\PY{n}{graph\PYZus{}fourier}\PY{p}{(}\PY{n}{short\PYZus{}sample\PYZus{}data}\PY{p}{,} \PY{n}{short\PYZus{}sample\PYZus{}fs}\PY{p}{,} \PY{n}{t}\PY{p}{,} \PY{n}{dt}\PY{p}{,} \PY{n}{title}\PY{p}{,} \PY{l+m+mi}{9}\PY{p}{,} \PY{l+m+mi}{5000}\PY{p}{,} \PY{n}{coef}\PY{o}{=}\PY{k+kc}{True}\PY{p}{,} \PY{n}{ticks}\PY{o}{=}\PY{l+m+mi}{250}\PY{p}{)}
\end{Verbatim}
\end{tcolorbox}

    \vspace{-0.50em}
    \begin{center}
    \adjustimage{max size={\linewidth}{\paperheight}}{tpe_files/tpe_36_0.png}
    \end{center}
    { \hspace*{\fill} \\}
    \vspace{-2.75em}
    
    Para la segunda vocal ``i'', se identifica en la señal el intervalo
\texttt{{[}0.087,\ 0.094{]}s} el cual corresponde con el periodo de una
señal periodica. Calculando los coeficientes de Fourier para dicho
período resulta como se muestra a continuación en la figura 10.

    \begin{tcolorbox}[breakable, size=fbox, boxrule=1pt, pad at break*=1mm,colback=cellbackground, colframe=cellborder, fontupper=\footnotesize]
\prompt{In}{incolor}{17}{\boxspacing}
\begin{Verbatim}[commandchars=\\\{\}]
\PY{n}{t}\PY{o}{=}\PY{l+m+mf}{0.087}
\PY{n}{dt}\PY{o}{=}\PY{l+m+mf}{0.007}
\PY{n}{title} \PY{o}{=} \PY{l+s+sa}{f}\PY{l+s+s1}{\PYZsq{}}\PY{l+s+s1}{Coeficientes de la Serie de Fourier para el intervalo [}\PY{l+s+si}{\PYZob{}}\PY{n}{t}\PY{l+s+si}{\PYZcb{}}\PY{l+s+s1}{, }\PY{l+s+si}{\PYZob{}}\PY{n}{t}\PY{o}{+}\PY{n}{dt}\PY{l+s+si}{\PYZcb{}}\PY{l+s+s1}{] del archivo `}\PY{l+s+si}{\PYZob{}}\PY{n}{short\PYZus{}sample\PYZus{}file\PYZus{}name}\PY{l+s+si}{\PYZcb{}}\PY{l+s+s1}{`}\PY{l+s+s1}{\PYZsq{}}
\PY{n}{graph\PYZus{}fourier}\PY{p}{(}\PY{n}{short\PYZus{}sample\PYZus{}data}\PY{p}{,} \PY{n}{short\PYZus{}sample\PYZus{}fs}\PY{p}{,} \PY{n}{t}\PY{p}{,} \PY{n}{dt}\PY{p}{,} \PY{n}{title}\PY{p}{,} \PY{l+m+mi}{10}\PY{p}{,} \PY{l+m+mi}{5000}\PY{p}{,} \PY{n}{coef}\PY{o}{=}\PY{k+kc}{True}\PY{p}{,} \PY{n}{ticks}\PY{o}{=}\PY{l+m+mi}{250}\PY{p}{)}
\end{Verbatim}
\end{tcolorbox}

    \vspace{-0.50em}
    \begin{center}
    \adjustimage{max size={\linewidth}{\paperheight}}{tpe_files/tpe_38_0.png}
    \end{center}
    { \hspace*{\fill} \\}
    \vspace{-2.75em}
    
    Tomando ahora 5 periodos de la señal, por ejemplo el intervalo entre
\texttt{{[}0.087,\ 0.122{]}}, los coeficientes de la serie de Fourier
resultan como se muestra en la figura 11.

    \begin{tcolorbox}[breakable, size=fbox, boxrule=1pt, pad at break*=1mm,colback=cellbackground, colframe=cellborder, fontupper=\footnotesize]
\prompt{In}{incolor}{18}{\boxspacing}
\begin{Verbatim}[commandchars=\\\{\}]
\PY{n}{t}\PY{o}{=}\PY{l+m+mf}{0.087}
\PY{n}{dt}\PY{o}{=}\PY{l+m+mi}{5}\PY{o}{*}\PY{l+m+mf}{0.007}
\PY{n}{title} \PY{o}{=} \PY{l+s+sa}{f}\PY{l+s+s1}{\PYZsq{}}\PY{l+s+s1}{Coeficientes de la Serie de Fourier para el intervalo [}\PY{l+s+si}{\PYZob{}}\PY{n}{t}\PY{l+s+si}{\PYZcb{}}\PY{l+s+s1}{, }\PY{l+s+si}{\PYZob{}}\PY{n}{t}\PY{o}{+}\PY{n}{dt}\PY{l+s+si}{\PYZcb{}}\PY{l+s+s1}{] del archivo `}\PY{l+s+si}{\PYZob{}}\PY{n}{short\PYZus{}sample\PYZus{}file\PYZus{}name}\PY{l+s+si}{\PYZcb{}}\PY{l+s+s1}{`}\PY{l+s+s1}{\PYZsq{}}
\PY{n}{graph\PYZus{}fourier}\PY{p}{(}\PY{n}{short\PYZus{}sample\PYZus{}data}\PY{p}{,} \PY{n}{short\PYZus{}sample\PYZus{}fs}\PY{p}{,} \PY{n}{t}\PY{p}{,} \PY{n}{dt}\PY{p}{,} \PY{n}{title}\PY{p}{,} \PY{l+m+mi}{11}\PY{p}{,} \PY{l+m+mi}{5000}\PY{p}{,} \PY{n}{coef}\PY{o}{=}\PY{k+kc}{True}\PY{p}{,} \PY{n}{ticks}\PY{o}{=}\PY{l+m+mi}{250}\PY{p}{)}
\end{Verbatim}
\end{tcolorbox}

    \vspace{-0.50em}
    \begin{center}
    \adjustimage{max size={\linewidth}{\paperheight}}{tpe_files/tpe_40_0.png}
    \end{center}
    { \hspace*{\fill} \\}
    \vspace{-2.75em}
    
    Finalmente, para la vocal ``o'', se identifica el intervalo
\texttt{{[}0.634,\ 0.643{]}s}, graficando los coeficientes de Fourier,
resultan como en la figura 12 a continuación.

    \begin{tcolorbox}[breakable, size=fbox, boxrule=1pt, pad at break*=1mm,colback=cellbackground, colframe=cellborder, fontupper=\footnotesize]
\prompt{In}{incolor}{19}{\boxspacing}
\begin{Verbatim}[commandchars=\\\{\}]
\PY{n}{t}\PY{o}{=}\PY{l+m+mf}{0.634}
\PY{n}{dt}\PY{o}{=}\PY{l+m+mf}{0.009}
\PY{n}{title} \PY{o}{=} \PY{l+s+sa}{f}\PY{l+s+s1}{\PYZsq{}}\PY{l+s+s1}{Coeficientes de la Serie de Fourier para el intervalo [}\PY{l+s+si}{\PYZob{}}\PY{n}{t}\PY{l+s+si}{\PYZcb{}}\PY{l+s+s1}{, }\PY{l+s+si}{\PYZob{}}\PY{n}{t}\PY{o}{+}\PY{n}{dt}\PY{l+s+si}{\PYZcb{}}\PY{l+s+s1}{] del archivo `}\PY{l+s+si}{\PYZob{}}\PY{n}{short\PYZus{}sample\PYZus{}file\PYZus{}name}\PY{l+s+si}{\PYZcb{}}\PY{l+s+s1}{`}\PY{l+s+s1}{\PYZsq{}}
\PY{n}{graph\PYZus{}fourier}\PY{p}{(}\PY{n}{short\PYZus{}sample\PYZus{}data}\PY{p}{,} \PY{n}{short\PYZus{}sample\PYZus{}fs}\PY{p}{,} \PY{n}{t}\PY{p}{,} \PY{n}{dt}\PY{p}{,} \PY{n}{title}\PY{p}{,} \PY{l+m+mi}{12}\PY{p}{,} \PY{l+m+mi}{5000}\PY{p}{,} \PY{n}{coef}\PY{o}{=}\PY{k+kc}{True}\PY{p}{,} \PY{n}{ticks}\PY{o}{=}\PY{l+m+mi}{250}\PY{p}{)}
\end{Verbatim}
\end{tcolorbox}

    \vspace{-0.50em}
    \begin{center}
    \adjustimage{max size={\linewidth}{\paperheight}}{tpe_files/tpe_42_0.png}
    \end{center}
    { \hspace*{\fill} \\}
    \vspace{-2.75em}
    
    Y tomando mas periodos, por ejemplo para aproximadamente 6 periodos, se
toma el intervalo \texttt{{[}0.634,\ 0.588{]}s}, el grafico de los
coeficientes de Fourier resulta como se muestra en la figura 13.

    \begin{tcolorbox}[breakable, size=fbox, boxrule=1pt, pad at break*=1mm,colback=cellbackground, colframe=cellborder, fontupper=\footnotesize]
\prompt{In}{incolor}{20}{\boxspacing}
\begin{Verbatim}[commandchars=\\\{\}]
\PY{n}{t}\PY{o}{=}\PY{l+m+mf}{0.634}
\PY{n}{dt}\PY{o}{=}\PY{l+m+mi}{6}\PY{o}{*}\PY{l+m+mf}{0.009}
\PY{n}{title} \PY{o}{=} \PY{l+s+sa}{f}\PY{l+s+s1}{\PYZsq{}}\PY{l+s+s1}{Coeficientes de la Serie de Fourier para el intervalo [}\PY{l+s+si}{\PYZob{}}\PY{n}{t}\PY{l+s+si}{\PYZcb{}}\PY{l+s+s1}{, }\PY{l+s+si}{\PYZob{}}\PY{n}{t}\PY{o}{+}\PY{n}{dt}\PY{l+s+si}{\PYZcb{}}\PY{l+s+s1}{] del archivo `}\PY{l+s+si}{\PYZob{}}\PY{n}{short\PYZus{}sample\PYZus{}file\PYZus{}name}\PY{l+s+si}{\PYZcb{}}\PY{l+s+s1}{`}\PY{l+s+s1}{\PYZsq{}}
\PY{n}{graph\PYZus{}fourier}\PY{p}{(}\PY{n}{short\PYZus{}sample\PYZus{}data}\PY{p}{,} \PY{n}{short\PYZus{}sample\PYZus{}fs}\PY{p}{,} \PY{n}{t}\PY{p}{,} \PY{n}{dt}\PY{p}{,} \PY{n}{title}\PY{p}{,} \PY{l+m+mi}{13}\PY{p}{,} \PY{l+m+mi}{2000}\PY{p}{,} \PY{n}{coef}\PY{o}{=}\PY{k+kc}{True}\PY{p}{,} \PY{n}{ticks}\PY{o}{=}\PY{l+m+mi}{200}\PY{p}{)}
\end{Verbatim}
\end{tcolorbox}

    \vspace{-0.50em}
    \begin{center}
    \adjustimage{max size={\linewidth}{\paperheight}}{tpe_files/tpe_44_0.png}
    \end{center}
    { \hspace*{\fill} \\}
    \vspace{-2.75em}
    

    % Add a bibliography block to the postdoc
    
    
    
\end{document}
